\chapter{Reproducibility Details}
\label{app:reproducibility}
\section{Source revision and environment}

Paths in this appendix are relative to the repository root. The thesis
was completed against the detector and documentation tree at revision
\texttt{16b33d6}, with the thesis changes subsequently applied. The
detector code preceding the first writing installment is identified by
\texttt{63ba80d}. These writing revisions do not replace the experiment
identities recorded in the individual bundles.

The closing final-partition protocol records source revision
\texttt{99072f0}. Its 109 dependency fingerprints cover code,
configurations, starting pairs, processed counts, labels, and original
annotations. Three Python source files subsequently changed; their
current hashes are verified against
\path{reports/scoring_order_verification.json}. The original frozen
hashes remain in the closing protocol and are copied alongside the
current ones in the thesis evidence manifest.

The historical resource image has a different boundary. Its environment
record identifies base revision \texttt{8365870}, a dirty working tree,
and the exact package-source and image hashes. Its image identity is
\begin{quote}\footnotesize
\path{sha256:76c00897e51f521264a4b2a87d0e963e3d330a7c4bb515c5e3ae4e42d41863b5}.
\end{quote}
The numerical records identify Python~3.12.14, NumPy~2.5.1, and
SciPy~1.18.0. The resource image identifies Debian~13 and
Python~3.11.17. Neither environment should be inferred from the
machine used only to compile this document.

\section{Evidence map}
\label{app:evidence-map}

Table~\ref{tab:evidence-map} maps the main results to their machine-readable
records. The generated file
\path{thesis/generated/evidence_manifest.json} enumerates presentation
source SHA-256 values, the frozen dependency count, and the three
superseded source records. A source hash provides identity, not an
assessment of scientific quality.

\begin{longtable}{@{}>{\raggedright\arraybackslash}p{.29\textwidth}
 >{\raggedright\arraybackslash}p{.66\textwidth}@{}}
\caption{Retained evidence used by the thesis. Directory names identify
the bundle; JSON files provide exact settings and source boundaries.}
\label{tab:evidence-map}\\
\toprule
Claim or result & Principal source paths \\
\midrule
\endfirsthead
\toprule
Claim or result & Principal source paths \\
\midrule
\endhead
\bottomrule
\endfoot
Reference partitions & \path{configs/d_link_day_cam5.json},
\path{configs/unsw_chromecast.json},
\path{configs/unsw_samsung_camera.json}; matching
\path{data/processed/} and \path{data/labels/} device directories. \\
Capture preparation & Device-specific \path{capture_inventory.json},
\path{capture_validation.json}, and \path{capture_statistics.json}
where retained under \path{reports/}; UNSW preparation records under
\path{unsw_iot_attack_pcaps/}. \\
Normal models and calibration & \path{artifacts/<device>/model.json}
and \path{artifacts/<device>/threshold.json}. \\
Reference static detection &
\path{reports/<device>/evaluation_statistics.json}; corresponding
\path{artifacts/<device>/events.json}. \\
Historical adaptive results &
\path{reports/experiments/legacy/<device>/<variant>/}; matching
\path{configs/experiments/legacy/} and
\path{artifacts/experiments/legacy/}. \\
Joint-scoring selection & \path{reports/detection_decision.json} and
\path{reports/detection_rejected_experiments.json}. \\
Backend agreement and timing &
\path{reports/D-LinkDayCam5/lm_backend_benchmark.json} and
\path{reports/detection_backend_checks.json}. \\
Native setup optimisation &
\path{docs/lm_backend_change_record.tex}, including the historical
paired timing and 1,000-matrix comparison. \\
Aggregate numerical experiment &
\path{reports/D-LinkDayCam5/lm_professor_experiment.json}. \\
Matched durations & \path{reports/window_size/results.json},
\path{reports/window_size/decisions.csv}, and
\path{artifacts/experiments/window_size/}. \\
Closing replay & \path{reports/final_evaluation/protocol.json},
\path{results.json}, and \path{verification.json} in that directory;
complete starting pairs, events, and attempts under
\path{artifacts/final_evaluation/}. \\
Later scoring regression &
\path{reports/scoring_order_verification.json}. \\
Resource measurements & \path{reports/edge/simulation-2000.json},
\path{reports/edge/simulation-20000.json}, and
\path{reports/edge/simulation-environment.json}. \\
Current thesis verification & \path{out/thesis/project_verification.log},
document build log, and final document-validation records. These are
generated outputs rather than replacements for historical experiments. \\
\end{longtable}

\section{Regenerating the document}

From the repository root, the following command verifies retained
presentation inputs, regenerates table rows and plot data, and runs
LuaLaTeX/Biber through the existing Docker Compose service:
\begin{verbatim}
./thesis/build.sh
\end{verbatim}
The output is \path{out/thesis/thesis.pdf}. Auxiliary files stay in
the same ignored output directory. The document uses the locally
available \path{texlive/texlive:latest} image with networking disabled.
A complete environment record and image identity are retained with
the final validation output; the mutable image tag alone is not
a permanent version identifier.

To regenerate only the presentation evidence, use
\begin{verbatim}
python3 thesis/generate_evidence.py
\end{verbatim}
This standard-library script reads the retained data and reports;
it does not preprocess packet captures, refit models, or replay
detectors. Assertions check confusion-count arithmetic and expected
dependencies. Generated TeX rows and CSV plot data remain under
\path{thesis/generated/}. They are presentation artifacts and must
not be edited to alter an experimental result.

The dissertation uses the cited subset of the annotated shortlist in
\path{docs/thesis_reference_review.tex}. That review distinguishes
full-text inspection from metadata-only verification. The bibliography
file \path{thesis/references.bib} holds publication identifiers and
links; existence of a DOI alone is not evidence for a particular
claim. The author should consult the original cited passages when
reviewing the final argument.

\section{Repeating detector experiments}
\label{app:experiment-commands}

\subsection{Installation and reference stages}

The following commands describe experiment entry points; they were
not executed merely to regenerate the thesis. Research scripts can
write their configured output paths. Use an isolated checkout and
copies of retained artifacts when repeating a historical experiment,
and retain its original revision before comparing new results.

A local installation requires Python and the native build tools.
The repository's documented native build is
\begin{verbatim}
gcc -O3 -fPIC -shared \
  src/lm_idnet/algorithms/native/LM_time_lib.c \
  -lm -o src/lm_idnet/algorithms/native/LM_time_lib.so
python -m pip install -e '.[test]'
\end{verbatim}
The Dockerfile provides the packaged alternative. Source packet
captures are required for preprocessing, but are not required to
score the retained processed windows. Their access and redistribution
remain subject to the publisher's terms and available downloads.

For a chosen device configuration, the reference stage order is
\begin{verbatim}
lm-idnet preprocess --config configs/unsw_chromecast.json
lm-idnet diagnose --config configs/unsw_chromecast.json
lm-idnet train --config configs/unsw_chromecast.json
lm-idnet calibrate --config configs/unsw_chromecast.json
lm-idnet score --config configs/unsw_chromecast.json
lm-idnet evaluate --config configs/unsw_chromecast.json
\end{verbatim}
Diagnosis and preparation outputs should be reviewed before replacing
an accepted reference. Repeating \texttt{train} or \texttt{calibrate}
changes the starting pair and therefore defines a new comparison
unless the resulting bytes match the retained identity.

\subsection{Adaptive, duration, and closing comparisons}

The conservative and historical development suites have separate
configuration and output directories:
\begin{verbatim}
./scripts/run_adaptation_experiments.sh
./scripts/run_adaptation_experiments.sh legacy
\end{verbatim}
These launch through the project runner. An individual adaptive
configuration can instead be supplied to \texttt{lm-idnet adapt},
followed by \texttt{lm-idnet evaluate}. Omitting a persistent state
directory starts from the configured reference pair; specifying
\texttt{--state-dir} resumes accepted state and changes that
experimental boundary.

The retained selected-method replay is generated by
\begin{verbatim}
python scripts/evaluate_final.py
\end{verbatim}
It writes its fixed bundle under \path{reports/final_evaluation/}
and \path{artifacts/final_evaluation/}. Running it on the current
source produces current count-first event scores. It should not
be expected to reproduce every historical full-score byte, even
when all binary decisions agree.

The duration script allows new output locations, for example
\begin{verbatim}
python scripts/evaluate_window_sizes.py \
  --output out/repeat-window-size/reports \
  --artifacts out/repeat-window-size/artifacts
\end{verbatim}
It prepares fresh model/calibration pairs at each duration and
retains decisions on the common support. Its separate fit/calibration
roles explain small differences from reference ten-minute metrics.
The aggregate numerical script is
\begin{verbatim}
python scripts/lm_professor_experiment.py
\end{verbatim}
Its fixed output path is
\path{reports/D-LinkDayCam5/lm_professor_experiment.json}.
Backend production timing is exposed through the \texttt{benchmark}
CLI; \texttt{lm-idnet benchmark --help} identifies its current
configuration and repetition arguments. Timing results should be
compared under the same environment and call boundary.

\section{Runtime and verification commands}

\path{docs/edge-deployment.md} supplies the full prepared-state,
offline-container, and isolated resource-measurement commands.
The short launcher input is
\begin{verbatim}
./scripts/run-edge.sh < tests/data/edge/windows.jsonl
\end{verbatim}
It requires a prebuilt image and prepared writable state. The
five-record fixture emits four events. The longer measurement
uses \path{scripts/benchmark_edge.py}, a fresh state copy, explicit
CPU/memory limits, and either 2,000 or 20,000 input records.
The repeated synthetic workload should not be used to estimate
real-traffic detection precision or recall.

The full current verification command used during thesis completion is
\begin{verbatim}
LM_IDNET_EDGE_IMAGE=lm-idnet:final \
  .venv/bin/python -m pytest -q -m ""
\end{verbatim}
It reports 643 passes without skips. The empty marker expression
includes the slow and benchmark-marked cases. The locally available
image supplies the opt-in container checks; this successful run is
not evidence that a fresh Docker Hub installation was performed.
The historical clean image build was blocked by registry DNS, and
the retained source-refreshed image boundary remains explicit.

The final PDF is separately checked for unresolved citations and
cross-references, layout overflow, and PDF/A conformance. These
document checks validate the rendered artifact rather than its
scientific conclusions. The final validation record gives the
actual tool versions, file identity, page count, and outcome.
